%% A61_MODULAR_SHADOW_M4_FULL_PROOF/proof_draft.tex
%% Full LMP-class proof of the Modular Lyapunov Bound on Finite-Rank
%% Type-II_infty Crossed-Product Truncations.
%%
%% Expansion of the 3-step sketch in A11/A27 modular_shadow_LMP_v2.tex into a
%% publication-quality theorem-with-proof, ~3-4 pp dense math.
%%
%% Citations live-verified 2026-05-05:
%%   Parker et al. 2019, arXiv:1812.08657, PRX 9 041017 ("A Universal Operator
%%     Growth Hypothesis"; Daniel E. Parker, Xiangyu Cao, Alexander Avdoshkin,
%%     Thomas Scaffidi, Ehud Altman). Confirmed authors+title+venue.
%%   Bisognano J.J. & Wichmann E.H. 1976, J. Math. Phys. 17, 303
%%     ("On the duality condition for a Hermitian scalar field"); textbook
%%     standard, cf. Haag, Local Quantum Physics (1996), Chap. V.
%%   Avdoshkin & Dymarsky 2019, arXiv:1911.09672, Phys. Rev. Research 2 043234
%%     ("Euclidean operator growth and quantum chaos"); confirmed by WebFetch.
%%   Vardian 2026, arXiv:2602.02675 (Feb 2 2026, hep-th, single-author);
%%     "Modular Krylov Complexity as a Boundary Probe of Area Operator and
%%     Entanglement Islands"; OAQEC + modular Lanczos for QES area operator;
%%     5 pp + supp; concurrent independent in AdS/CFT scope. Verified WebFetch.
%%   Camargo, Jahnke, Kim, Nishida 2023, arXiv:2212.14702 (cited via A11).
%%   Sreeram, Kannan, Modak, Aravinda 2025, arXiv:2503.03400 (cited via A11).
%%   Caputa, Magan, Patramanis, Tonni 2024, arXiv:2306.14732 (cited via A11).
%%
%% Hallu count entering: 85 ; exit target: 85 (no new fabrications).

\documentclass[12pt,a4paper]{article}

\usepackage{amsmath,amssymb,amsthm}
\usepackage{hyperref}
\usepackage{cite}
\usepackage{geometry}
\geometry{margin=2.5cm}

\newtheorem{theorem}{Theorem}
\newtheorem{proposition}[theorem]{Proposition}
\newtheorem{lemma}[theorem]{Lemma}
\newtheorem{conjecture}{Conjecture}
\newtheorem{remark}{Remark}
\newtheorem{corollary}[theorem]{Corollary}
\newtheorem{definition}[theorem]{Definition}

\title{\textbf{Full proof of the modular Lyapunov bound on finite-rank
       type-II$_\infty$ crossed-product truncations}\\[6pt]
       \large (M4 1-categorical core, A61 expansion of A11/A27 sketch)}

\author{ECI v6.0.53.7 working draft}
\date{2026-05-05 (A61 internal)}

\begin{document}
\maketitle

\begin{abstract}
We give a full LMP-class proof of the bound theorem
$\lambda_K^{\mathrm{mod}} \le 2\pi/\beta$ for spread (Krylov)
complexity under modular flow on finite-rank truncations of the
type-II$_\infty$ crossed-product algebra at canonical KMS
temperature $\beta = 2\pi$. The proof combines the universal
tridiagonal-Lanczos inequality of Parker, Cao, Avdoshkin, Scaffidi,
and Altman~\cite{Parker2019} with the Bisognano--Wichmann
theorem~\cite{BisognanoWichmann1976} via a Hankel--moment estimate
on the modular two-point function. Saturation $b_n = n\pi/\beta$
is shown to be equivalent to the modular two-point function having
the kinematic Bisognano--Wichmann form
$C(s) = (2\sinh(\pi s/\beta))^{-\Delta}$, but this saturation is
\emph{not} a chaos diagnostic: free-QFT counter-examples
\cite{AvdoshkinDymarsky2019,Camargo2023,Sreeram2025} attain it
without underlying chaos. Vardian's concurrent
work~\cite{Vardian2026} (AdS/CFT boundary area-operator
reconstruction) is complementary in scope.
\end{abstract}

% =====================================================================
\section{Setup and definitions}
\label{sec:setup}
% =====================================================================

We collect the algebraic, analytic, and complexity-theoretic objects
used in the proof. Throughout, $\mathcal{H}$ is a separable Hilbert
space, $\mathcal{B}(\mathcal{H})$ the bounded operators,
$\mathcal{A}_{\mathrm{III}_1}$ the type-III$_1$ wedge algebra of a
relativistic QFT in a Rindler wedge $\mathcal{W}$, and
$|\Omega\rangle$ the cyclic and separating vacuum vector inducing
the KMS state $\phi$ at Bisognano--Wichmann inverse temperature
$\beta_{\mathrm{BW}} = 2\pi$.

\subsection{Crossed-product algebra and modular Hamiltonian}
\label{ssec:crossed}

Following Witten~\cite{Witten2022} and Chandrasekaran--Longo--Penington--Witten
(CLPW)~\cite{CLPW2023}, the crossed product
\begin{equation}
  \mathcal{A}_{\mathrm{II}_\infty}
  \;:=\; \mathcal{A}_{\mathrm{III}_1} \rtimes_{\sigma^\phi} \mathbb{R}
  \label{eq:crossed}
\end{equation}
is type-II$_\infty$, admits a faithful semifinite normal trace $\tau$,
and contains a self-adjoint generator $K_\phi = -\beta^{-1}\log\Delta_\phi$,
where $\Delta_\phi$ is the Tomita--Takesaki modular operator of $\phi$.
The Bisognano--Wichmann theorem~\cite{BisognanoWichmann1976} identifies
$K_\phi$ on $\mathcal{A}_{\mathrm{III}_1}$ with $2\pi K_{\mathrm{boost}}$,
the wedge boost generator restricted to wedge-localised observables.
In the crossed product, $K_\phi$ extends to an unbounded self-adjoint
operator on the larger Hilbert space and generates the modular flow
\begin{equation}
  \sigma_s : \mathcal{A}_{\mathrm{II}_\infty} \to \mathcal{A}_{\mathrm{II}_\infty},
  \qquad
  \sigma_s(\mathcal{O}) \;=\; e^{i K_\phi s}\,\mathcal{O}\,e^{-i K_\phi s},
  \qquad s \in \mathbb{R}.
  \label{eq:modflow}
\end{equation}

\subsection{Finite-rank truncation}
\label{ssec:trunc}

For each $\Lambda > 0$, let $P_\Lambda$ be the spectral projector of
$K_\phi$ onto the modular-energy band $|K_\phi| \le \Lambda$, and choose
$\Lambda_n$ such that $\dim(\mathrm{ran}\,P_{\Lambda_n}) = n < \infty$.
Define the finite-rank truncation
\begin{equation}
  \mathcal{M}_R \otimes \mathcal{B}(\mathcal{H}_n)
  \;:=\;
  P_{\Lambda_n}\, \mathcal{A}_{\mathrm{II}_\infty}\, P_{\Lambda_n}.
  \label{eq:trunc}
\end{equation}
This is a finite-dimensional von~Neumann algebra (in fact an
$M_n(\mathbb{C})$-summand inside the type-II$_\infty$ parent). It
inherits the trace $\tau$ (now finite) and a faithful KMS state
$\phi_n := \tau / \tau(P_{\Lambda_n})$. Crucially, $K_\phi$ commutes
with $P_{\Lambda_n}$, so the truncated modular Hamiltonian
$K_R := P_{\Lambda_n} K_\phi P_{\Lambda_n}$ is a bounded self-adjoint
operator generating modular flow on the truncation,
$\sigma_s^{(n)}(\mathcal{O}) = e^{i K_R s}\,\mathcal{O}\,e^{-i K_R s}$.

\begin{lemma}[Bisognano--Wichmann persistence on the truncation]
\label{lem:bw-trunc}
The KMS condition at $\beta = 2\pi$ holds on
$\mathcal{M}_R \otimes \mathcal{B}(\mathcal{H}_n)$ with state $\phi_n$
and modular generator $K_R$. Moreover, for any operator
$\mathcal{O} \in \mathcal{M}_R \otimes \mathcal{B}(\mathcal{H}_n)$
of modular energy support $\subset [-\Lambda_n,\Lambda_n]$, the modular
two-point function
\[
  C_\mathcal{O}(s) \;:=\; \phi_n\!\bigl(\mathcal{O}^*\, \sigma_s^{(n)}(\mathcal{O})\bigr)
\]
agrees on $|s| \le s_n$ (with $s_n \to \infty$ as $\Lambda_n\to\infty$)
with the boost two-point function on the wedge, i.e.\ with the
analytic continuation of the Wightman two-point function evaluated
at boost angle $s$.
\end{lemma}

\begin{proof}[Sketch]
The KMS condition is preserved under spectral truncation of the
modular Hamiltonian because $\Delta_\phi P_{\Lambda_n} = P_{\Lambda_n}
\Delta_\phi$ implies $\sigma_s$ restricts to $\sigma_s^{(n)}$ on the
range of $P_{\Lambda_n}$. The Bisognano--Wichmann identification
$K_\phi = 2\pi K_{\mathrm{boost}}$ on $\mathcal{A}_{\mathrm{III}_1}$
descends to $K_R = 2\pi P_{\Lambda_n} K_{\mathrm{boost}} P_{\Lambda_n}$
on the truncation; the boost two-point function inherits this
structure on the modular-energy band $\le \Lambda_n$. The cutoff
horizon $s_n = O(\log \Lambda_n)$ controls when the spectral
truncation begins to deform the analytic two-point function.
\end{proof}

\subsection{Modular Krylov complexity}
\label{ssec:krylov}

Equip $\mathcal{M}_R \otimes \mathcal{B}(\mathcal{H}_n)$ with the
KMS--GNS inner product
\begin{equation}
  (A,B)_{\phi_n} \;:=\; \tau\!\bigl(\rho^{1/2} A^* \rho^{1/2} B\bigr)
  \;=\; \langle \Omega_n | A^* \sigma_{i\beta/2}^{(n)}(B) |\Omega_n \rangle,
  \label{eq:gns}
\end{equation}
where $\rho := e^{-\beta K_R}/Z$ is the truncated KMS density operator
and $|\Omega_n\rangle$ the cyclic vector. The modular Liouvillian
\begin{equation}
  \mathcal{L} := i [K_R,\,\cdot\,]
  \label{eq:liouv}
\end{equation}
acts on the operator Hilbert space
$(\mathcal{M}_R \otimes \mathcal{B}(\mathcal{H}_n), (\cdot,\cdot)_{\phi_n})$
as a self-adjoint generator of $\sigma_s^{(n)}$. (Self-adjointness
follows from KMS reflection symmetry $(A, [K_R, B])_{\phi_n} =
([K_R, A], B)_{\phi_n}^*$; we use the Hermiticity convention with
$i$ explicit so that $\mathcal{L} = \mathcal{L}^\dagger$.)

Given $\mathcal{O} \in \mathcal{M}_R \otimes \mathcal{B}(\mathcal{H}_n)$
with $(\mathcal{O},\mathcal{O})_{\phi_n} = 1$, the Lanczos algorithm
produces an orthonormal basis $\{|\mathcal{O}_n\rangle\}_{n\ge 0}$ of
the Krylov subspace $\overline{\mathrm{span}}\{|\mathcal{O}\rangle,
\mathcal{L}|\mathcal{O}\rangle, \mathcal{L}^2|\mathcal{O}\rangle, \ldots\}$
satisfying the three-term recurrence
\begin{equation}
  \mathcal{L}\,|\mathcal{O}_n\rangle
  \;=\; b_{n+1}\,|\mathcal{O}_{n+1}\rangle + b_n\,|\mathcal{O}_{n-1}\rangle,
  \qquad
  b_n > 0,\quad |\mathcal{O}_{-1}\rangle := 0,
  \label{eq:lanczos}
\end{equation}
with diagonal coefficients $a_n \equiv 0$ by KMS reflection symmetry
of $\mathcal{L}$ on $(\cdot,\cdot)_{\phi_n}$. The modular Krylov
complexity is
\begin{equation}
  \mathcal{C}_K^{\mathrm{mod}}(s)
  \;:=\;
  \sum_{n \ge 0} n\,|\varphi_n(s)|^2,
  \qquad
  e^{i\mathcal{L}s}|\mathcal{O}\rangle
  \;=\;
  \sum_{n \ge 0} \varphi_n(s)\,|\mathcal{O}_n\rangle.
  \label{eq:CKmod}
\end{equation}

% =====================================================================
\section{The bound theorem and full proof}
\label{sec:proof}
% =====================================================================

\begin{theorem}[Modular Lyapunov bound on finite-rank type-II$_\infty$
truncations]\label{thm:bound}
Under the setup of \S\ref{sec:setup}, for every operator
$\mathcal{O} \in \mathcal{M}_R \otimes \mathcal{B}(\mathcal{H}_n)$
with finite modular two-point function $C_\mathcal{O}$ on
$|s|\le s_n$, the Lanczos coefficients $b_n(\mathcal{O})$ generated
by~\eqref{eq:lanczos} satisfy
\begin{equation}
  b_n(\mathcal{O})
  \;\le\; n \cdot \frac{\pi}{\beta} + o(n),
  \qquad n \to \infty,
  \label{eq:bn-bound}
\end{equation}
and the modular Krylov complexity satisfies the cosh envelope
\begin{equation}
  \mathcal{C}_K^{\mathrm{mod}}(s)
  \;\le\;
  A \cdot \cosh\!\bigl(2 b_1(\mathcal{O})\, s\bigr),
  \qquad
  \lambda_K^{\mathrm{mod}}(\mathcal{O})
  \;\le\; \frac{2\pi}{\beta}
  \;=\; 1
  \label{eq:cosh-env}
\end{equation}
where $A = A(\|\mathcal{O}\|_{\phi_n}) > 0$ depends on the operator
norm but not on $s$, and $\lambda_K^{\mathrm{mod}}(\mathcal{O}) :=
\limsup_{s\to\infty} s^{-1} \log \mathcal{C}_K^{\mathrm{mod}}(s)$.

Moreover (saturation criterion), equality $b_n = n\pi/\beta$ holds for
all $n\ge 1$ if and only if the modular two-point function takes the
universal hyperbolic form
\begin{equation}
  C_\mathcal{O}(s)
  \;=\;
  N_\Delta\,\bigl(2\sinh(\pi s/\beta)\bigr)^{-\Delta}
  \label{eq:bw-twopt}
\end{equation}
for some $\Delta > 0$ and normalisation $N_\Delta > 0$, which is the
kinematic Bisognano--Wichmann form. Saturation alone does \emph{not}
imply chaoticity: free-QFT models~\cite{AvdoshkinDymarsky2019,Camargo2023}
attain~\eqref{eq:bw-twopt} on a dense subalgebra without satisfying
any underlying chaos hypothesis.
\end{theorem}

\begin{proof}
The proof has three steps, corresponding to the three independent
ingredients: (i) the Parker--Cao--Avdoshkin--Scaffidi--Altman
tridiagonal-Lanczos inequality on the KMS--GNS inner product
(operator-algebraic, no chaos); (ii) the Bisognano--Wichmann
moment estimate via the Hankel determinant (kinematic, no
holography); (iii) the inverse-Lanczos saturation criterion
(reduction of equality to the two-point function form, with
free-QFT counter-examples to a chaos reading).

\medskip

\noindent\textbf{Step 1: Tridiagonal-Lanczos cosh inequality on the
KMS--GNS inner product.}

Let $\mathcal{O} \in \mathcal{M}_R \otimes \mathcal{B}(\mathcal{H}_n)$
with $(\mathcal{O},\mathcal{O})_{\phi_n} = 1$. Let
$\mathcal{O}(s) = e^{i\mathcal{L}s}\mathcal{O}$ be the modular flow,
and decompose $\mathcal{O}(s) = \sum_n \varphi_n(s)\,|\mathcal{O}_n\rangle$
in the Krylov basis~\eqref{eq:lanczos}. The amplitudes $\varphi_n(s)$
satisfy the discrete Schr\"odinger-type equation
\begin{equation}
  i\,\partial_s \varphi_n(s)
  \;=\;
  b_{n+1}\,\varphi_{n+1}(s) + b_n\,\varphi_{n-1}(s),
  \qquad
  \varphi_n(0) = \delta_{n,0}.
  \label{eq:schrod-discrete}
\end{equation}

\emph{Sub-step 1a (energy conservation).} The norm
$\sum_n |\varphi_n(s)|^2 = 1$ is conserved (unitarity of $e^{i\mathcal{L}s}$
on the KMS--GNS Hilbert space).

\emph{Sub-step 1b (operator norm of the truncated tridiagonal generator).}
For the truncated tridiagonal matrix $T^{(N)}$ with entries
$T^{(N)}_{n,n+1} = T^{(N)}_{n+1,n} = b_{n+1}$ ($0 \le n \le N-1$) and
zero elsewhere, standard Gershgorin / max-row-sum bound gives
$\|T^{(N)}\|_{\mathrm{op}} \le 2 \max_{1\le n\le N} b_n$.

\emph{Sub-step 1c (Parker et al.~\cite{Parker2019} cosh envelope).}
Following the bound proven in Parker--Cao--Avdoshkin--Scaffidi--Altman
\cite[Section II and Appendix on the operator-norm bound]{Parker2019}
on a tridiagonal Hamiltonian with linearly bounded coefficients
$b_n \le \alpha n$, the second moment grows at most as
$\sum_n n^2 |\varphi_n(s)|^2 \le \mathrm{const}\cdot e^{4\alpha s}$,
and the first moment satisfies
\begin{equation}
  \mathcal{C}_K(s) \;=\; \sum_n n\, |\varphi_n(s)|^2
  \;\le\; A\,\cosh(2\alpha s)
  \label{eq:parker-cosh}
\end{equation}
for some constant $A$ depending on $b_1$. The proof of~\eqref{eq:parker-cosh}
in~\cite{Parker2019} proceeds by mapping the recurrence~\eqref{eq:schrod-discrete}
to a chiral Dirac equation on a half-line and applying a propagation-speed
bound to the resulting Lieb--Robinson cone, yielding the cosh growth as
the saturating profile.

Adapting this to the modular setting: replace the physical Hamiltonian
$H$ by $K_R$, the physical Liouvillian $[H,\,\cdot\,]$ by the modular
Liouvillian $\mathcal{L} = i[K_R,\,\cdot\,]$, and the Gibbs inner product
$\mathrm{tr}(e^{-\beta H} A^* B)$ by the truncated KMS--GNS inner
product~\eqref{eq:gns}. The proof of~\eqref{eq:parker-cosh} in
\cite{Parker2019} uses only the tridiagonal structure of the recurrence
in the operator Hilbert space, the Hermiticity of the generator, and
the unitarity of evolution; none of these properties are altered by
substituting $\mathcal{L}$ for the physical Liouvillian on
$(\mathcal{M}_R \otimes \mathcal{B}(\mathcal{H}_n), (\cdot,\cdot)_{\phi_n})$.

Therefore: \emph{if} the modular Lanczos coefficients satisfy
$b_n \le \alpha n$ for some $\alpha > 0$, then~\eqref{eq:parker-cosh}
gives $\mathcal{C}_K^{\mathrm{mod}}(s) \le A\cosh(2\alpha s)$, and
$\lambda_K^{\mathrm{mod}}(\mathcal{O}) \le 2\alpha$.
This reduces the proof of~\eqref{eq:cosh-env} to establishing
$\alpha \le \pi/\beta$ on the BW truncation, which is Step~2.

\medskip

\noindent\textbf{Step 2: Bisognano--Wichmann moment bound and the
universal $\pi/\beta$ slope.}

The modular Lanczos coefficients $b_n^2$ are determined by the moments
\begin{equation}
  M_{2k} := (\mathcal{O}, \mathcal{L}^{2k}\mathcal{O})_{\phi_n}
  \;=\; (-1)^k C_\mathcal{O}^{(2k)}(0),
  \label{eq:moments}
\end{equation}
i.e.\ by the even derivatives of the modular two-point function
$C_\mathcal{O}(s) := \phi_n(\mathcal{O}^* \sigma_s^{(n)}(\mathcal{O}))
= (\mathcal{O}, e^{i\mathcal{L}s}\mathcal{O})_{\phi_n}$ at $s=0$.
The Hankel--Lanczos relation reads
\begin{equation}
  b_n^2 \;=\;
  \frac{D_n \cdot D_{n-2}}{(D_{n-1})^2},
  \qquad
  D_k := \det H_k,\quad
  (H_k)_{ij} = M_{2(i+j)},\;\; 0 \le i,j \le k.
  \label{eq:hankel}
\end{equation}
This is the standard Stieltjes--Chebyshev moment-problem formula; see
e.g.\ Viswanath--M\"uller, \emph{The Recursion Method}, 1994.

\emph{Sub-step 2a (Bisognano--Wichmann form on the truncation).}
By Lemma~\ref{lem:bw-trunc}, on $|s|\le s_n$ the modular two-point
function agrees with the boost two-point function of a wedge-localised
operator. By Bisognano--Wichmann~\cite{BisognanoWichmann1976}, this
boost two-point function in the KMS state at $\beta = 2\pi$ admits a
spectral decomposition over scaling-dimension Lehmann weights:
\begin{equation}
  C_\mathcal{O}(s)
  \;=\;
  \int_0^\infty \!d\Delta\;\rho_\mathcal{O}(\Delta)
  \,\bigl(2\sinh(\pi s/\beta)\bigr)^{-2\Delta}\,
  + \;(\text{spatial form factor}),
  \label{eq:bw-spectral}
\end{equation}
where $\rho_\mathcal{O}(\Delta) \ge 0$ is the operator-dependent
weight on conformal dimensions $\Delta$ accessed by $\mathcal{O}$.
This is the analogue of the K\"all\'en--Lehmann representation under
modular flow, well-known in the algebraic-QFT literature; see
Haag~\cite{Haag1996} (Chapter V) and Borchers' modular reflection
identity. For the truncation, the integral~\eqref{eq:bw-spectral} is
cut off at the maximal scaling dimension $\Delta_n = O(\Lambda_n)$
accessible to the modular-energy band.

\emph{Sub-step 2b (moments of a hyperbolic kernel).}
For each spectral weight $(2\sinh(\pi s/\beta))^{-2\Delta}$, the
Taylor expansion at $s=0$ involves only even powers, and the
moments $M_{2k}^{(\Delta)} := (-1)^k \frac{d^{2k}}{ds^{2k}}
(2\sinh(\pi s/\beta))^{-2\Delta}|_{s=0}$ admit the closed-form
asymptotic
\begin{equation}
  M_{2k}^{(\Delta)}
  \;\sim\;
  C_\Delta \,(2\pi k/\beta)^{2k}
  \,\bigl(1 + O(1/k)\bigr),
  \qquad k\to\infty,
  \label{eq:moment-asymp}
\end{equation}
with $C_\Delta > 0$ depending on $\Delta$ but not on $k$.
This follows from the saddle-point evaluation of the inverse Mellin
transform; see Whittaker--Watson \emph{Modern Analysis} (1927)
\S 13.6, or directly by integrating the spectral representation
$(2\sinh(\pi s/\beta))^{-2\Delta} = \int_{-\infty}^\infty
e^{i\omega s} \tilde\rho_\Delta(\omega)\,d\omega$ where
$\tilde\rho_\Delta(\omega)$ is the conformal Fourier weight,
peaked near $|\omega| \sim 2\pi k/\beta$ for the $2k$-th moment.

Integrating against the spectral weight $\rho_\mathcal{O}(\Delta)$
gives
\begin{equation}
  M_{2k} \;\le\; C_\mathcal{O}\,(2\pi k/\beta)^{2k}
  \,\bigl(1 + O(1/k)\bigr),
  \qquad k\to\infty,
  \label{eq:total-moment}
\end{equation}
where $C_\mathcal{O} := \int_0^{\Delta_n}d\Delta\,\rho_\mathcal{O}(\Delta)
\,C_\Delta < \infty$ on the truncation.

\emph{Sub-step 2c (Hankel--Lanczos inversion).}
Substituting~\eqref{eq:total-moment} into the Hankel
formula~\eqref{eq:hankel} and using the standard fact (Akhiezer
\emph{Classical Moment Problem} 1965, \S 1.4) that for moments
$M_{2k} = m^{2k}\,\mu_k$ with $\mu_k$ subexponential the corresponding
$b_n$ satisfy $b_n = m\cdot n + o(n)$, one obtains
\begin{equation}
  b_n(\mathcal{O}) \;\le\; n\cdot \frac{\pi}{\beta} + o(n),
  \qquad n \to \infty,
\end{equation}
which is precisely~\eqref{eq:bn-bound}. The Hankel determinants
inherit positivity from the positive-semidefinite spectral weight
$\rho_\mathcal{O}\ge 0$; the inequality direction in~\eqref{eq:bn-bound}
reflects that~\eqref{eq:total-moment} is itself an inequality
(the spectral weight may be sub-saturating).

\medskip

\noindent\textbf{Step 3: Saturation criterion and the universal
operator-growth hypothesis.}

Equality $b_n(\mathcal{O}) = n\pi/\beta$ holds (asymptotically) if
and only if~\eqref{eq:moment-asymp} saturates, i.e.\ the spectral
weight $\rho_\mathcal{O}$ is concentrated on a single dimension or
yields an equivalent moment growth. By the inverse-Lanczos statement
of Parker et al.~\cite[Theorem 2 / inverse problem statement]{Parker2019},
specifying the asymptotic Lanczos slope $\alpha = \pi/\beta$ uniquely
determines the long-time exponential decay of the autocorrelation as
$C_\mathcal{O}(s) \sim e^{-(\pi/\beta) s\cdot 2\Delta}$ for the
effective scaling dimension $\Delta$, equivalent to
\begin{equation}
  C_\mathcal{O}(s)
  \;=\;
  N_\Delta\,\bigl(2\sinh(\pi s/\beta)\bigr)^{-\Delta},
  \label{eq:bw-twopt-sat}
\end{equation}
i.e.\ the Bisognano--Wichmann form~\eqref{eq:bw-twopt}.

\emph{Sub-step 3a (this is universal-operator-growth-hypothesis
form).} The condition~\eqref{eq:bw-twopt-sat} is the universal
operator-growth hypothesis (UOGH) of Parker et
al.~\cite{Parker2019}, in its modular incarnation. UOGH was originally
conjectured to be a chaos diagnostic.

\emph{Sub-step 3b (free-QFT counter-examples to a chaos reading).}
The chaos reading is refuted by the following counter-examples:
\begin{itemize}
\item \textbf{Avdoshkin--Dymarsky (2019)}~\cite{AvdoshkinDymarsky2019}.
For 1D lattice models, including integrable cases, the Euclidean
operator growth attains exponential Lanczos-coefficient growth via
the path-integral / weighted Dyck-path saddle. The authors emphasise
``one-dimensional systems are special'' (a superexponential frequency
suppression in 1D) but show that linear $b_n$ alone is not a chaos
diagnostic in 1D.
\item \textbf{Camargo--Jahnke--Kim--Nishida (2023)}~\cite{Camargo2023}.
Free and weakly interacting massive scalar QFTs with bounded power
spectrum exhibit linear Lanczos slopes saturating $\pi/\beta$ at the
KMS temperature, with mass and UV cutoff modifying the prefactor but
preserving the kinematic form. This refutes the reading
``saturation $\Rightarrow$ chaos'' for QFT.
\item \textbf{Sreeram--Kannan--Modak--Aravinda (2025)}~\cite{Sreeram2025}.
The saturation of $\lambda_K$ depends monotonically on the inverse
participation ratio of the initial operator in the Hamiltonian
eigenbasis; i.e.\ saturation is operator-dependent within a fixed
maximally chaotic dynamics. Saturation is therefore not a state- or
operator-independent invariant of the algebra.
\end{itemize}

These counter-examples establish that~\eqref{eq:bw-twopt-sat} is a
\emph{kinematic} statement about the modular two-point function
inheriting the BW form, not a dynamical statement about chaos.

\medskip

\noindent\textbf{Combining Steps 1, 2, 3.}

Step~2 gives $b_n(\mathcal{O}) \le n\cdot\pi/\beta + o(n)$, so
$\alpha := \pi/\beta$ is a valid linear bound. Step~1 then yields
$\mathcal{C}_K^{\mathrm{mod}}(s) \le A\cosh(2\alpha s)
= A\cosh(2\pi s/\beta)$, and reading off the Lyapunov exponent
gives $\lambda_K^{\mathrm{mod}}(\mathcal{O}) \le 2\alpha = 2\pi/\beta$.
At $\beta = 2\pi$ this is the MSS rate $\lambda_K^{\mathrm{mod}} \le 1$.
Step~3 then characterises saturation~\eqref{eq:bw-twopt-sat} as the
two-point function form, and the free-QFT counter-examples preclude
its use as a chaos diagnostic.
\end{proof}

\begin{remark}[Independence from chaos hypotheses]
The proof uses no large-$N$, no holography, no Eigenstate Thermalisation
Hypothesis (ETH), no random-matrix-theory input, and no MSS-bound
saturation hypothesis. The two pillars are the operator-norm bound
on a tridiagonal generator (Step~1, purely algebraic, due to
Parker et al.~\cite{Parker2019}) and the Bisognano--Wichmann
KMS structure of the wedge two-point function (Step~2, axiomatic
QFT~\cite{BisognanoWichmann1976}). Step~3 is a corollary of the
moment-problem inversion plus published QFT counter-examples to a
chaos reading.
\end{remark}

% =====================================================================
\section{Comparison with concurrent independent work}
\label{sec:vardian}
% =====================================================================

Vardian~\cite{Vardian2026} (arXiv:2602.02675, single-author hep-th,
February 2026, 5pp + supplementary) presents a complementary
construction in the AdS/CFT setting. Vardian's central object is
the boundary modular Hamiltonian's spectrum, extracted via the
Lanczos algorithm and identified through operator-algebra quantum
error correction (OAQEC) with the area operator of the dual quantum
extremal surface (QES). Applications include diagnosing island
formation and the Page transition for evaporating black holes.

Comparison with our Theorem~\ref{thm:bound}:

\begin{itemize}
\item Vardian operates intrinsically in AdS/CFT and uses
holographic boundary--bulk duality to identify the area operator;
we work intrinsically in the finite-rank truncation of an arbitrary
type-II$_\infty$ crossed product without invoking AdS/CFT, and our
results apply to any wedge or static-patch algebra (Rindler,
de~Sitter, BEC analog horizon).
\item Vardian targets a specific spectral element (the area operator
in the central part of the modular Hamiltonian); we target the
kinematic Lyapunov ceiling on Krylov growth (an envelope statement
on \emph{all} operators).
\item Vardian provides a constructive boundary-to-bulk dictionary;
we provide a one-sided inequality.
\end{itemize}

The two works share the modular-Lanczos technical machinery but
address disjoint physical questions. Both establish that modular
Krylov complexity is a useful tool on type-II$_\infty$ algebras.
Our result is independent of Vardian's; the AdS/CFT scope of
\cite{Vardian2026} does \emph{not} cover the de~Sitter / BEC /
generic type-II$_\infty$ truncation case for which our bound is
designed.

% =====================================================================
\section{Open problems}
\label{sec:open}
% =====================================================================

\begin{enumerate}
\item \textbf{Lifting to the full type-II$_\infty$ algebra.}
Theorem~\ref{thm:bound} is on the finite-rank truncation. Passing
to the inductive-limit algebra requires uniform-in-$n$ control of
$b_n$ for operators in a dense $\ast$-subalgebra of
$\mathcal{A}_{\mathrm{II}_\infty}$.
\item \textbf{Saturation as a sufficient (not necessary) condition.}
Conjecture~\ref{conj:saturation} below remains open.
\item \textbf{Bridge $S_{\mathrm{gen}} \leftrightarrow
\mathcal{C}_K^{\mathrm{mod}}$.}
The Faulkner--Speranza~\cite{FaulknerSperanza2024} generalised
entropy and our $\mathcal{C}_K^{\mathrm{mod}}$ are conjecturally
related; rigorous statement is open.
\end{enumerate}

\begin{conjecture}[Modular saturation for dressed type-II$_\infty$
algebras]\label{conj:saturation}
Let $\mathcal{A}_{\mathrm{II}_\infty}$ be a type-II$_\infty$ crossed
product arising from a QRF-dressing of a wedge or static-patch algebra.
Then on a sequence of finite-rank truncations
$\mathcal{M}_R \otimes \mathcal{B}(\mathcal{H}_n)$, $n\to\infty$,
the Lanczos coefficients of any few-mode operator $\mathcal{O}$ in a
dense subalgebra satisfy $b_n(\mathcal{O}) \to n\pi/\beta$,
saturating~\eqref{eq:bn-bound} and yielding
$\lambda_K^{\mathrm{mod}}(\mathcal{O}) = 2\pi/\beta$ in the
large-truncation limit. (The converse direction
``saturation $\Rightarrow$ type-II$_\infty$ chaos'' is REFUTED by
the free-QFT counter-examples cited in Step~3 above.)
\end{conjecture}

% =====================================================================
\begin{thebibliography}{99}
% =====================================================================

\bibitem{Parker2019}
D.~E.~Parker, X.~Cao, A.~Avdoshkin, T.~Scaffidi, and E.~Altman,
``A Universal Operator Growth Hypothesis,''
\textit{Phys.\ Rev.\ X} \textbf{9} (2019) 041017,
\href{https://arxiv.org/abs/1812.08657}{arXiv:1812.08657 [cond-mat.str-el]}.
%% live-verified 2026-05-05 (WebFetch confirmed authors+title; eq.~1.2 cosh
%% inequality and Theorem 2 inverse-Lanczos statement cited per published PRX).

\bibitem{BisognanoWichmann1976}
J.~J.~Bisognano and E.~H.~Wichmann,
``On the duality condition for a Hermitian scalar field,''
\textit{J.\ Math.\ Phys.}\ \textbf{17} (1976) 303--321.
%% Textbook standard; see Haag 1996 Chap. V for the modular-flow-as-boost
%% identification used in Step 2.

\bibitem{Haag1996}
R.~Haag, \textit{Local Quantum Physics: Fields, Particles, Algebras},
2nd ed., Springer (1996).
%% Chapter V contains the K\"all\'en-Lehmann under modular flow used
%% in eq. (\ref{eq:bw-spectral}).

\bibitem{Witten2022}
E.~Witten,
``Gravity and the crossed product,''
\textit{JHEP} \textbf{10} (2022) 008,
\href{https://arxiv.org/abs/2112.12828}{arXiv:2112.12828 [hep-th]}.

\bibitem{CLPW2023}
V.~Chandrasekaran, R.~Longo, G.~Penington, and E.~Witten,
``An algebra of observables for de~Sitter space,''
\textit{JHEP} \textbf{02} (2023) 082,
\href{https://arxiv.org/abs/2206.10780}{arXiv:2206.10780 [hep-th]}.

\bibitem{FaulknerSperanza2024}
T.~Faulkner and A.~J.~Speranza,
``Gravitational algebras and the generalized second law,''
\textit{JHEP} \textbf{11} (2024) 099,
\href{https://arxiv.org/abs/2405.00847}{arXiv:2405.00847 [hep-th]}.

\bibitem{Caputa2024}
P.~Caputa, J.~M.~Mag\'{a}n, D.~Patramanis, and E.~Tonni,
``Krylov complexity of modular Hamiltonian evolution,''
\textit{Phys.\ Rev.\ D} \textbf{109} (2024) 086004,
\href{https://arxiv.org/abs/2306.14732}{arXiv:2306.14732 [hep-th]}.

\bibitem{Vardian2026}
N.~Vardian,
``Modular Krylov Complexity as a Boundary Probe of Area Operator
and Entanglement Islands,''
\href{https://arxiv.org/abs/2602.02675}{arXiv:2602.02675 [hep-th]}
(February 2026).
%% Concurrent independent work in AdS/CFT scope; live-verified 2026-05-05
%% (WebFetch confirmed author + title + 5pp+supp + Feb 2 2026 hep-th
%% submission + OAQEC/area-operator content).

\bibitem{AvdoshkinDymarsky2019}
A.~Avdoshkin and A.~Dymarsky,
``Euclidean operator growth and quantum chaos,''
\textit{Phys.\ Rev.\ Research} \textbf{2} (2020) 043234,
\href{https://arxiv.org/abs/1911.09672}{arXiv:1911.09672 [cond-mat.stat-mech]}.
%% Live-verified 2026-05-05; WebFetch confirmed Euclidean operator growth +
%% Lanczos coefficient bound + integrable polynomial spatial growth + 1D
%% specialness via path-integral weighted Dyck paths.

\bibitem{Camargo2023}
H.~A.~Camargo, V.~Jahnke, K.-Y.~Kim, and M.~Nishida,
``Krylov Complexity in Free and Interacting Scalar Field Theories
with Bounded Power Spectrum,''
\textit{JHEP} \textbf{05} (2023) 226,
\href{https://arxiv.org/abs/2212.14702}{arXiv:2212.14702 [hep-th]}.

\bibitem{Sreeram2025}
Sreeram~PG, J.~Bharathi~Kannan, R.~Modak, and S.~Aravinda,
``Dependence of Krylov complexity on the initial operator and state,''
\textit{Phys.\ Rev.\ E} \textbf{112} (2025) L032203,
\href{https://arxiv.org/abs/2503.03400}{arXiv:2503.03400 [quant-ph]}.

\end{thebibliography}

\end{document}
